\documentclass[10pt,a4paper]{amsart}
\usepackage[margin=19mm]{geometry}
\usepackage{amsmath,amssymb,mathtools}
\usepackage[scr=rsfs,cal=euler]{mathalfa}
\usepackage{xfrac}
\usepackage[unicode,hidelinks,bookmarks=false]{hyperref}
\newcommand{\ie}{i.e.\ }
\newcommand{\cf}{cf.\ }
\newcommand{\resp}{respectively\ }
\newcommand{\Subsection}{\subsection}
\providecommand{\nobreakdash}{\nobreak\hbox{-}}
\setlength{\emergencystretch}{2em}
\multlinegap0pt
\usepackage{bm}
\usepackage{bbm}
\usepackage{dsfont}
\usepackage{xspace}
\usepackage{cancel}
\usepackage{mathtools}
\usepackage{amsthm}
\makeatletter
\@ifundefined{lemma}{\newtheorem{lemma}{Lemma}}{}
\makeatother
\newcommand{\ad}{\mathrm{ad}}
\title[Ricci curvature for hydrodynamics on the sphere]{Ricci curvature for hydrodynamics on the sphere}
\author{Leandro Lichtenfelz, Klas Modin, Stephen C. Preston}
\date{Source: arXiv:2508.09833v1 (CC BY 4.0)}
\begin{document}
\maketitle

\noindent\emph{Source and licence.} Ricci curvature for hydrodynamics on the sphere. Version arXiv:2508.09833v1 (CC BY 4.0), distributed under the Creative Commons Attribution 4.0 International licence, as recorded in the arXiv licence metadata for this exact version and shown on the abstract page https://arxiv.org/abs/2508.09833v1. Full rights notice: licenses/arxiv-2508.09833-CC-BY-4.0.txt.\par\smallskip

\noindent\textit{Editorial scope.} Counted unit: the Lemma whose source label is \texttt{lemma\_A\_XY} in the article (source0.tex lines 1027--1034, source label \texttt{lemma\_A\_XY}), delivered together with the article's own complete original proof of it (source0.tex lines 1036--1085, one \texttt{proof} environment of 50 lines). No mathematics is written, repaired or re-derived: statement and proof are the article's own text line for line. Required context retained uncounted: eq\_zeitlin\_metric (lines 212--216); eq\_coefficients (lines 502--516); lemma\_3\_sum (lines 937--945). Every definition, display and label of those lines is retained unchanged. Omitted from this package and not redistributed: 1--211, 217--501, 517--936, 946--1026, 1035--1035, 1086--1202. Nothing inside a retained excerpt is omitted or altered: every retained body is delivered exactly as the article's lines are written, so no omission receipt of any kind accompanies this package. Citations: the retained excerpts carry no citation command at all, so no bibliography entry of the article is transcribed and no reference is redistributed. Reference adaptation: none was needed and none was made; every reference inside the delivered unit resolves inside it, so no unresolved reference or citation is produced. Printed slips kept literal: no printed word, symbol, hypothesis or number of the counted statement or of its proof was altered. Layout and class: the article's own class is \texttt{amsart} with packages including amsmath, amsfonts, amssymb, mathrsfs, amsthm, graphicx, xcolor, float, diagbox, todonotes, ulem, hyperref; the wrapper is the lane's pinned \texttt{amsart} editorial wrapper at 10pt on A4, which restates the theorem-like environments the retained text needs and adds the article's own macro definitions that the retained text needs (\texttt{\textbackslash ad}), with extra wrapper packages (bm, bbm, dsfont, xspace, cancel, mathtools). The delivered numbering is the unit's own. Assets: no retained span contains a figure environment, a graphics-inclusion command, a PDF-page inclusion, a table or a TikZ picture; no image file is copied into this package, the package declares no asset dependency and no image file is redistributed with it. Licence: version arXiv:2508.09833v1 of this article is distributed under the Creative Commons Attribution 4.0 International licence, as recorded in the arXiv licence metadata for this exact version and shown on the abstract page https://arxiv.org/abs/2508.09833v1. A full rights notice is delivered at licenses/arxiv-2508.09833-CC-BY-4.0.txt. Characters: every retained excerpt is pure ASCII, so the delivered engine, XeLaTeX with its default Latin Modern text font, reports no missing character.
\begin{align}
(u, v) &= \mathrm{Tr}\big(u^\dagger v\big), \qquad &\text{(bi-invariant metric)} \label{eq_biinvariant_metric}
\\
\langle u, v\rangle &= \frac{1}{N} \operatorname{Tr}\big(u^\dagger (-\Delta_N) v\big), &\text{(Zeitlin metric)} \label{eq_zeitlin_metric}
\end{align}

\begin{gather}\label{eq_coefficients}
\begin{split}
C_{\ell m\ell' m'}^{\ell'' m''(N)} &= S_{\ell,\ell',\ell''} \Big(1 - (-1)^{\ell + \ell' + \ell''}\Big)(-1)^{m''+1}
\left(\begin{matrix}
\ell & \ell' & \ell'' \\
m & m' & m''
\end{matrix}
\right)
\left\{\begin{matrix}
\ell & \ell' & \ell'' \\
\frac{N-1}{2} & \frac{N-1}{2} & \frac{N-1}{2}
\end{matrix}
\right\},
\end{split}
\end{gather}

\begin{lemma}\label{lemma_3_sum}
Let $C_{\ell,m,\ell',m'}^{\ell'',m''}$ denote the structure constants as given in \eqref{eq_coefficients}. Then, we have
\begin{equation}
\Big(C_{\ell,m,\ell,1-m}^{1,-1}\Big)^2 + 
\Big(C_{\ell,m,\ell,-m}^{1,0}\Big)^2 + 
\Big(C_{\ell,m,\ell,-m-1}^{1,1}\Big)^2 = \frac{12\ell(\ell+1)}{N(N^2-1)}
\end{equation}    
for all $1 \leq \ell \leq N-1$ and $|m| \leq \ell$. Note that the right-hand side does not depend on $m$.
\end{lemma}

\begin{lemma}\label{lemma_A_XY}
For any fixed $(\ell,m)$ with $2 \leq \ell \leq N-1$ and $-\ell \leq m \leq \ell$, we have
\begin{align*}
\sum\limits_{\ell' = 2}^{N-1}
~\sum\limits_{m' = -\ell'}^{\ell'} 3\frac{\|A_{T_{\ell m}}T_{\ell'm'}\|^2}{\|T_{\ell m} \wedge T_{\ell'm'}\|^2} = \frac{9}{2\ell(\ell+1)},
\end{align*}
where in the above summation, we skip the term $(\ell', m') = (\ell, m)$, for which the summand on the left is not well-defined.
\end{lemma}

\begin{proof}
We expand the Lie bracket of $X = T_{\ell m}$ and $Y = T_{\ell'm'}$ according to  \eqref{eq_coefficients}, but since we are only interested in the vertical component, the condition $\ell'' = 1$ is imposed, which forces  $m''\in\{-1,0,1\}$, and we get 
\begin{gather}\label{eq_A_formula_1}
\begin{split}
A_{T_{\ell m}}T_{\ell' m'} &= \frac{1}{2}\mathcal{V} (\ad_{T_{\ell m}} T_{\ell'm'})  =-\frac{1}{2\hbar_N} 
[T_{\ell m}, T_{\ell' m'}] \\
&= -\frac{1}{2\hbar_N}\Big(C_{\ell m\ell' m'}^{1,-1}T_{1,-1} + C_{\ell m\ell' m'}^{1,0}T_{1,0} + C_{\ell m\ell' m'}^{1,1}T_{1,1}\Big).
\end{split}
\end{gather}
Selection rules for $3j$ symbols imply a dependence of $\ell', m'$ on $\ell, m$, so the only nonzero values are:
\begin{gather}\label{eq_A_formula_2}
\begin{split}
A_{T_{\ell m}}T_{\ell, 1-m}
&=
-\tfrac{1}{2\hbar_N}C_{\ell,m,\ell,1-m}^{1,-1}T_{1,-1} \qquad \text{for $m \neq -\ell$}, \\
A_{T_{\ell m}}T_{\ell, -m} 
&=
-\tfrac{1}{2\hbar_N}C_{\ell,m,\ell,-m}^{1,0}T_{1,0}, \\
A_{T_{\ell m}}T_{\ell, -1-m}
&=
-\tfrac{1}{2\hbar_N}C_{\ell,m,\ell,-m-1}^{1,1}T_{1,1} \qquad \text{for $m \neq \ell$}.
\end{split}
\end{gather}
Recall that the $\{ T_{\ell m} \}$ basis is orthonormal in the bi-invariant metric, but only orthogonal in the Zeitlin metric \eqref{eq_zeitlin_metric}, so we have \begin{equation*}
\| T_{\ell m} \wedge T_{\ell' m'} \|^2 = \|T_{\ell m}\|^2\|T_{\ell'm'}\|^2 = \frac{\ell(\ell+1)\ell'(\ell'+1)}{N^2}, \qquad \text{for $(\ell, m) \neq (\ell',m')$}.
\end{equation*}
In particular, $\|T_{1m}\|^2 = 2/N$ for $|m| \leq 1$. Thus, summing over all $(\ell', m') \neq (\ell, m)$ and using \eqref{eq_A_formula_2},
\begin{gather*}
\begin{split}
&\sum\limits_{\ell' = 2}^{N-1}
\sum\limits_{m' = -\ell'}^{\ell'} 3\frac{\|A_{T_{\ell m}}T_{\ell'm'}\|^2}{\|T_{\ell m} \wedge T_{\ell'm'}\|^2}
=
\frac{3N^2}{\ell^2(\ell+1)^2} \sum\limits_{\ell' = 2}^{N-1}
\sum\limits_{m' = -\ell'}^{\ell'} \|A_{T_{\ell m}}T_{\ell'm'}\|^2 \\
&\qquad\qquad=
\frac{3N^2}{4\ell^2(\ell+1)^2\hbar_N^2}\left(
\Big\|C_{\ell,m,\ell,1-m}^{1, -1}T_{1, -1}\Big\|^2 +
\Big\|C_{\ell,m,\ell,1-m}^{1, 0}T_{1, 0}\Big\|^2 +
\Big\|C_{\ell,m,\ell,1-m}^{1, 1}T_{1, 1}\Big\|^2\right) \\
&\qquad\qquad= 
\frac{3N}{2\ell^2(\ell+1)^2\hbar_N^2}\left(
\left(C_{\ell,m,\ell,1-m}^{1, -1}\right)^2 +
\left(C_{\ell,m,\ell,1-m}^{1, 0}\right)^2 +
\left(C_{\ell,m,\ell,1-m}^{1, 1}\right)^2\right) \\
&\qquad\qquad= \frac{3N}{2\ell^2(\ell+1)^2}\frac{N^2-1}{4}\frac{12\ell(\ell+1)}{N(N^2-1)} \\
&\qquad\qquad= \frac{9}{2\ell(\ell+1)},
\end{split}
\end{gather*}
by Lemma \ref{lemma_3_sum}.
\end{proof}

\end{document}
